% ============================================================
\section{Failure Mode Analysis}
\label{app:failure_mode}
% ============================================================

While \system\ consistently outperforms baselines
across all configurations, it is not infallible.
We identify three failure patterns that arise in
practice.  Importantly, none stems from a design
flaw in the calibration mechanism itself; rather,
each is triggered by data-level conditions that
degrade the quality of the signals $\Phi$ relies on.

\subsection{Sparse Memory at Early Rounds}
\label{app:failure_sparse}

$\Phi$'s scoring passes evaluate $\Delta_t$ against
the existing memory state $\mathcal{M}_t$.  When
$\mathcal{M}_t$ contains very few entries (typically
during the first $1$--$2$ debate rounds) the
signals become unreliable.

\begin{tcolorbox}[colback=gray!5, colframe=gray!60,
  title={\small Example: HotpotQA, MemBank, Round 1},
  fonttitle=\bfseries\small, boxrule=0.4pt,
  arc=1.5mm]
\small
\textbf{Task}: \emph{``Which film came first,
Interstellar or The Martian?''}\\[3pt]
\textbf{$\Delta_t$}: ``Interstellar was released in
2014. The Martian was released in 2015.
Therefore Interstellar came first.''
\hfill\textit{(correct)}\\[3pt]
\textbf{$\rho_{\text{detect}}$}: Memory contains
only 2 prior entries, both unrelated to these films.
$\hat{r}_k = 0$ for all new vectors $\Rightarrow$
$\rho_{\text{detect}} = 1.0$ (vacuously high).\\[3pt]
\textbf{$\rho_{\text{align}}$}: Auditor queries
return empty top-$k$ lists from the near-empty
memory, so $d_{\text{RBO}} = 0.0$ for all auditors
$\Rightarrow$ $\rho_{\text{align}} = 1.0$
(vacuously high).\\[3pt]
\textbf{Outcome}: $\rho = 1.0$.  The entry is
committed with full trust, which happens to be
correct here.  However, if $\Delta_t$ had been
incorrect, $\Phi$ would have assigned the same score,
because there is no existing evidence to calibrate
against.
\end{tcolorbox}

\noindent
\textbf{Analysis.}
This is not a flaw in $\Phi$'s design but a
data-availability issue: calibration requires a
sufficiently populated $\mathcal{M}_t$ to produce
meaningful scores.  The adaptive threshold $\rho^*$
(running median of recent scores) partially mitigates
this (early rounds produce uniformly high $\rho$
values), so $\rho^*$ is also high, and only entries
that are \emph{relatively} worse than their peers
are filtered.  Nevertheless, the absolute quality of
early-round gating is weaker than later rounds.
Figure~\ref{fig:debate_rounds} confirms this: the
performance gap between \system\ and Vanilla is
smallest at round~1 and widens as $\mathcal{M}_t$
grows.

\subsection{Correlated Hallucinations Across Agents}
\label{app:failure_correlated}

$\rho_{\text{align}}$ measures whether $\Delta_t$
disrupts retrieval relative to auditor queries.
This signal loses power when multiple agents share
the same misconception, because the auditors'
queries are themselves biased toward the error.

\begin{tcolorbox}[colback=gray!5, colframe=gray!60,
  title={\small Example: LoCoMo, G-Memory, Round 4},
  fonttitle=\bfseries\small, boxrule=0.4pt,
  arc=1.5mm]
\small
\textbf{Context}: A long conversation mentions that
``Alice moved to Boston in March'' (session 3) and
``Alice started a new job in April'' (session 5).
The ground truth is that the job is in Boston.\\[3pt]
\textbf{$\Delta_t$}: ``Alice's new job is in
New York.'' \hfill\textit{(incorrect)}\\[3pt]
\textbf{What happens}: The backbone LLM
(Qwen3-VL-8B-Instruct) has a parametric bias toward
associating ``new job'' with ``New York'' (a frequent
co-occurrence in pre-training data).  Four out of
five auditors have the same bias in their reasoning
context, so their intercepted queries are about
``Alice's job in New York'' rather than ``Alice's
job in Boston.''\\[3pt]
\textbf{$\rho_{\text{align}}$}: Since $\Delta_t$
is consistent with the majority of auditor queries,
retrieval disruption is low $\Rightarrow$
$\rho_{\text{align}} = 0.72$.\\[3pt]
\textbf{$\rho_{\text{detect}}$}: The embedding of
$\Delta_t$ does not cluster unusually densely
$\Rightarrow$ $\rho_{\text{detect}} = 0.68$.\\[3pt]
\textbf{Outcome}: $\rho = \sqrt{0.72 \times 0.68}
= 0.70 > \rho^* = 0.55$.  \textbf{Committed
(incorrectly).}
\end{tcolorbox}

\noindent
\textbf{Analysis.}
When agents share a correlated hallucination (often
caused by the backbone LLM's parametric biases) the
``crowd'' of auditor probes reinforces rather than
challenges the error.  $\Phi$ is designed to detect
manipulation by a \emph{minority} of agents
(Section~\ref{sec:problem}), not systematic biases
shared by the majority.  This failure mode is
fundamentally a \emph{data quality} issue: the
auditor probes are only as good as the diversity of
the agents' reasoning.  Using heterogeneous backbones
(e.g., mixing Qwen with Llama or Mistral) would
reduce probe correlation, as different LLMs have
different parametric biases.  We note that this
failure mode also affects all baselines equally: LLM Audit and PPL Filter are even more susceptible,
since they evaluate each entry independently without
any crowd signal at all.

\subsection{Ambiguous Graph Semantics}
\label{app:failure_graph}

The graph calibration passes ($s_{\text{local}}$,
$s_{\text{path}}$) rely on the relation types and
entity nodes extracted by the memory architecture.
When the extraction produces ambiguous or
inconsistent relation labels, calibration quality
degrades.

\begin{tcolorbox}[colback=gray!5, colframe=gray!60,
  title={\small Example: ALFWorld, AriGraph, Round 3},
  fonttitle=\bfseries\small, boxrule=0.4pt,
  arc=1.5mm]
\small
\textbf{Existing edges in $\mathcal{M}_t$}:\\
\texttt{(apple\_1, located\_in, fridge\_1)}\\
\texttt{(fridge\_1, contains, apple\_1)}\\
\texttt{(apple\_1, state, cool)}\\[3pt]
\textbf{$\Delta_t$}: Agent extracts the edge
\texttt{(apple\_1, inside, fridge\_1)}.
\hfill\textit{(correct but redundant)}\\[3pt]
\textbf{$s_{\text{local}}$}: The relation
\texttt{inside} does not appear in
$\mathcal{R}_{\text{compat}}(\texttt{located\_in})$
because the co-occurrence statistics have not yet
observed \texttt{inside} and \texttt{located\_in}
together $\Rightarrow$ $s_{\text{local}} = 0.08$.\\[3pt]
\textbf{$s_{\text{path}}$}: BFS from
\texttt{apple\_1} to \texttt{fridge\_1} finds the
existing path via \texttt{located\_in}, but
$\text{comp}(\texttt{inside},
\texttt{located\_in}) = 0.15 < \eta_t = 0.40$
$\Rightarrow$ $s_{\text{path}} = 0.0$.\\[3pt]
\textbf{Outcome}: $\rho_{\text{detect}} = 0.03$,
leading to $\rho = 0.08 < \rho^*$.
\textbf{Rejected (false positive).}
\end{tcolorbox}

\noindent
\textbf{Analysis.}
The edge \texttt{(apple\_1, inside, fridge\_1)} is
semantically correct (\texttt{inside} and
\texttt{located\_in} are near-synonyms) but $\Phi$
rejects it because the compatibility relation
$\mathcal{R}_{\text{compat}}$ is estimated from
co-occurrence statistics in the graph, not from
external semantic knowledge.  Early in the episode,
when the graph is small, co-occurrence statistics
are noisy and may not capture synonym relations.
This is a data sparsity issue at the graph level,
analogous to the embedding sparsity issue in
\S\ref{app:failure_sparse}.

This failure mode has limited practical impact for
two reasons.  First, the rejected entry is
\emph{redundant} (the same fact is already
represented by \texttt{located\_in}) so rejecting
it does not lose information.  Second, as the graph
grows, $\mathcal{R}_{\text{compat}}$ accumulates
more co-occurrence evidence and eventually learns
that \texttt{inside} and \texttt{located\_in} are
compatible, at which point similar edges would pass.
The adaptive threshold $\eta_t$ (running median)
also adjusts downward as more diverse relation types
appear.

\subsection{Summary}

Table~\ref{tab:failure_summary} categorises the
three failure modes by their root cause and
mitigation.

\begin{table}[h]
\centering
\caption{Summary of failure modes.  All three are
  caused by data-level conditions, not by design
  limitations of $\Phi$.}
\label{tab:failure_summary}
\small
\setlength{\tabcolsep}{3pt}
\begin{tabular}{p{2.2cm}|p{2.5cm}|p{2.8cm}}
\toprule
\textbf{Failure mode} & \textbf{Root cause} &
\textbf{Mitigation} \\
\midrule
Sparse memory (early rounds)
  & Insufficient entries in $\mathcal{M}_t$ for
    meaningful calibration
  & Adaptive $\rho^*$ provides relative filtering;
    impact diminishes as memory grows \\
\midrule
Correlated hallucinations
  & Backbone LLM's parametric biases shared across
    agents
  & Heterogeneous backbones reduce probe correlation;
    affects all methods equally \\
\midrule
Ambiguous graph semantics
  & Noisy relation extraction from small graphs
  & $\mathcal{R}_{\text{compat}}$ self-corrects as
    graph grows; rejected entries are typically
    redundant \\
\bottomrule
\end{tabular}
\end{table}

\noindent
Across all three failure modes, we observe a common
pattern: $\Phi$'s calibration quality is bounded by
the quality and diversity of the data in
$\mathcal{M}_t$ and the agents' reasoning contexts.
As memory accumulates and agent interactions
diversify, all three failure modes naturally
attenuate.  This is consistent with the debate-round
analysis (Figure~\ref{fig:debate_rounds}), which
shows \system's advantage growing over time as
$\mathcal{M}_t$ becomes richer.